\paragraph{虚轴投影的正交通道与对偶逆平方门}
命题 \ref{prop:jouwkowsky-inverse-square} 把单位圆上的实投影 \(W(z)=z+z^{-1}\) 与逆平方门 \(J(z^2)\) 组合为 Lee--Yang 负实射线像。
在同一代数框架内，单位圆上的虚轴投影提供一个正交通道：其像仍由逆平方门给出，但落在正实射线上。

\begin{proposition}[虚轴投影门与正实射线像]\label{prop:leyang-orthogonal-dual-imaginary-projection}
\leanverified{paper\_leyang\_orthogonal\_dual\_imaginary\_projection}
定义
\[
W_{-}(z):=z-z^{-1}\qquad(z\in\CC^{\times}),
\qquad
J_{-}(u):=-\frac{u}{(1-u)^2}\qquad(u\in\CC\setminus\{1\}).
\]
则对任意 \(z\in\CC^{\times}\setminus\{\pm 1\}\) 成立严格恒等式
\[
J_{-}(z^{2})=-\frac{1}{W_{-}(z)^2}.
\]
若 \(|z|=1\)，则 \(W_{-}(z)\in \mathrm{i}[-2,2]\)，从而
\[
J_{-}(z^{2})\in\left[\frac14,\infty\right).
\]
并且对任意 \(\lambda>\frac14\)，方程 \(J_{-}(e^{2\mathrm{i}\theta})=\lambda\)（\(\theta\in[0,2\pi)\)）恰有四个解；在端点 \(\lambda=\frac14\) 处退化为两个解。
\end{proposition}
\begin{proof}
由恒等式
\[
\frac{1-z^{2}}{z}=z^{-1}-z=-W_{-}(z)
\]
得 \((1-z^{2})^{2}=z^{2}W_{-}(z)^{2}\)，故
\[
J_{-}(z^{2})
=-\frac{z^{2}}{(1-z^{2})^{2}}
=-\frac{1}{W_{-}(z)^{2}}.
\]
当 \(|z|=1\) 时写 \(z=e^{\mathrm{i}\theta}\)，则 \(W_{-}(z)=2\mathrm{i}\sin\theta\in \mathrm{i}[-2,2]\)，从而
\[
J_{-}(z^{2})=-\frac{1}{(2\mathrm{i}\sin\theta)^{2}}=\frac{1}{4\sin^{2}\theta}\in\left[\frac14,\infty\right),
\]
其中 \(\sin\theta=0\) 对应 \(z=\pm 1\) 的发散端点。
对任意 \(\lambda>\frac14\)，方程 \(\frac{1}{4\sin^{2}\theta}=\lambda\) 等价于 \(\sin^{2}\theta=\frac{1}{4\lambda}\in(0,1)\)，在 \([0,2\pi)\) 上具有四个互异解 \(\theta\mapsto\pm\theta\) 与 \(\theta\mapsto\pi\pm\theta\)；
在 \(\lambda=\frac14\) 时对应 \(\sin^{2}\theta=1\)，仅有 \(\theta=\frac{\pi}{2},\frac{3\pi}{2}\) 两解。证毕。
\end{proof}

\endinput
